% ============================================================================
% Penjelasan fundamental MLP + contoh numerik (untuk skripsi).
% Cara pakai: % ============================================================================
% Penjelasan fundamental MLP + contoh numerik (untuk skripsi).
% Cara pakai: % ============================================================================
% Penjelasan fundamental MLP + contoh numerik (untuk skripsi).
% Cara pakai: % ============================================================================
% Penjelasan fundamental MLP + contoh numerik (untuk skripsi).
% Cara pakai: \input{bab/penjelasan_mlp} di Bab 2 (teori) / Bab 3 (perancangan).
% Catatan: desimal di dalam math memakai titik (.) agar spasi rapi; teks biasa
%          boleh memakai koma. Butuh paket amsmath (umumnya sudah dimuat).
% ============================================================================

\section{Dasar Jaringan Saraf Tiruan (Multilayer Perceptron)}
\label{sec:dasar-mlp}

Pengendali pada penelitian ini menggunakan jaringan saraf tiruan jenis
\emph{Multilayer Perceptron} (MLP), yaitu jaringan \emph{feedforward} di mana data
mengalir satu arah dari lapisan masukan, melalui lapisan tersembunyi, menuju lapisan
keluaran. MLP terdiri atas unit-unit hitung (neuron) yang saling terhubung melalui
\emph{bobot} (\emph{weight}) dan disertai \emph{bias}. Pada tiap neuron, masukan
berbobot dijumlahkan lalu dilewatkan ke sebuah \emph{fungsi aktivasi} non-linear.

Arsitektur yang dirancang adalah \textbf{3--4--1}, yaitu 3 neuron masukan, 4 neuron
pada satu lapisan tersembunyi, dan 1 neuron keluaran:
\begin{itemize}
  \item \textbf{Masukan} ($\mathbf{x}$): galat tekanan $e$, perubahan galat $\Delta e$,
        dan \emph{duty cycle} saat ini $u$.
  \item \textbf{Keluaran} ($y$): perubahan \emph{duty cycle} $\Delta u$ ($\Delta Duty$).
\end{itemize}

\subsection{Fungsi Aktivasi}
\label{sub:aktivasi}
Lapisan tersembunyi memakai fungsi aktivasi \emph{hyperbolic tangent} (tanh):
\begin{equation}
\tanh(z) = \frac{e^{z}-e^{-z}}{e^{z}+e^{-z}}, \qquad \tanh(z)\in(-1,1).
\label{eq:tanh}
\end{equation}
Fungsi non-linear ini penting: tanpa aktivasi non-linear, gabungan beberapa lapisan
linear akan setara dengan satu fungsi linear saja, sehingga jaringan tidak dapat
memodelkan hubungan melengkung. Lapisan keluaran memakai aktivasi linear (identitas)
agar nilai $\Delta u$ bebas bernilai positif, nol, maupun negatif.

\subsection{Propagasi Maju (\emph{Forward Propagation})}
\label{sub:forward}
Misalkan $\mathbf{x'}$ adalah vektor masukan (setelah dinormalisasi). Keluaran
lapisan tersembunyi dan lapisan keluaran dihitung sebagai:
\begin{align}
\mathbf{z} &= \mathbf{x'}\,W_1 + \mathbf{b}_1, &
\mathbf{h} &= \tanh(\mathbf{z}), \label{eq:hidden}\\
y &= \mathbf{h}\,W_2 + b_2, & \Delta u &= y. \label{eq:output}
\end{align}
dengan $W_1$ matriks bobot lapisan 1, $\mathbf{b}_1$ vektor bias lapisan 1, $W_2$
matriks bobot lapisan 2, dan $b_2$ bias keluaran. Persamaan~\eqref{eq:hidden} dan
\eqref{eq:output} adalah inti perhitungan MLP: perkalian matriks, penjumlahan bias,
dan fungsi aktivasi.

\subsection{Normalisasi Masukan}
\label{sub:scaling}
Karena ketiga masukan memiliki rentang sangat berbeda (misal $e\approx0{,}1$ tetapi
$u\approx80$), tiap masukan distandardisasi terlebih dahulu menggunakan rata-rata
$\mu_i$ dan simpangan baku $\sigma_i$ dari data latih:
\begin{equation}
x'_i = \frac{x_i - \mu_i}{\sigma_i}.
\label{eq:scaling}
\end{equation}

\section{Perancangan dan Model Matematis Pengendali MLP}
\label{sec:perancangan-mlp}

\subsection{Bobot Terlatih dan Dimensi Matriks}
\label{sub:bobot}
Setelah pelatihan, diperoleh nilai bobot dan bias yang tetap. Untuk arsitektur
3--4--1:
\begin{equation}
W_1 =
\begin{bmatrix}
-1.1576 & 1.7171 & -3.6581 & -16.9100\\
0.0077 & 0.1740 & -3.6078 & 7.4493\\
0.9206 & 0.7704 & -11.9742 & 10.4903
\end{bmatrix},\quad
\mathbf{b}_1 =
\begin{bmatrix}
0.5008\\ -0.3952\\ -22.8359\\ -7.2012
\end{bmatrix}^{\!\top}
\label{eq:W1}
\end{equation}
\begin{equation}
W_2 =
\begin{bmatrix}
-1.3646\\ 1.7042\\ 0.9956\\ -0.3038
\end{bmatrix},\qquad
b_2 = 1.2567 .
\label{eq:W2}
\end{equation}
Dimensi matriks ditentukan oleh aturan perkalian matriks. Karena 3 masukan menuju 4
neuron, $W_1$ berukuran $3\times4$ (tiap baris satu fitur, tiap kolom satu neuron).
Karena 4 neuron menuju 1 keluaran, $W_2$ berukuran $4\times1$. Dengan demikian:
\begin{equation}
\underbrace{(1\times3)}_{\mathbf{x'}}\cdot\underbrace{(3\times4)}_{W_1}=(1\times4),
\qquad
\underbrace{(1\times4)}_{\mathbf{h}}\cdot\underbrace{(4\times1)}_{W_2}=(1\times1),
\end{equation}
yaitu dimensi dalam harus sama ($3=3$ dan $4=4$). Jumlah parameter total adalah
$(3\times4+4)+(4\times1+1)=21$.

\subsection{Contoh Perhitungan Numerik (Data Nyata)}
\label{sub:contoh}
Diambil satu kondisi nyata: $P^{*}=0{,}300$~bar, $P=0{,}2408$~bar, sehingga
$e=0{,}0592$; $\Delta e=-0{,}0047$; dan $u=82{,}00\%$. Maka
$\mathbf{x}=[\,0.0592,\,-0.0047,\,82.00\,]$.

\textbf{Langkah 1 --- Normalisasi} (Pers.~\ref{eq:scaling}), dengan
$\mu=[-0.0159,\,0.0010,\,82.272]$ dan $\sigma=[0.1159,\,0.0616,\,7.1318]$:
\begin{equation}
\mathbf{x'} = [\,0.6475,\ -0.0931,\ -0.0382\,].
\end{equation}

\textbf{Langkah 2 --- Pra-aktivasi lapisan tersembunyi} (Pers.~\ref{eq:hidden}).
Sebagai contoh, neuron pertama (kolom ke-1 dari $W_1$):
\begin{equation}
z_1 = 0.5008 + (0.6475)(-1.1576) + (-0.0931)(0.0077) + (-0.0382)(0.9206) = -0.2845.
\end{equation}
Untuk keempat neuron diperoleh:
\begin{equation}
\mathbf{z} = [\,-0.2845,\ 0.6710,\ -24.4117,\ -19.2442\,].
\end{equation}

\textbf{Langkah 3 --- Aktivasi tanh} (Pers.~\ref{eq:tanh}):
\begin{equation}
\mathbf{h} = \tanh(\mathbf{z}) = [\,-0.2771,\ 0.5857,\ -1.0000,\ -1.0000\,].
\end{equation}
(Dua neuron terakhir jenuh pada $-1$ karena nilai $z$ sangat negatif.)

\textbf{Langkah 4 --- Lapisan keluaran} (Pers.~\ref{eq:output}):
\begin{align}
y &= 1.2567 + (-0.2771)(-1.3646) + (0.5857)(1.7042) \notag\\
  &\quad + (-1.0000)(0.9956) + (-1.0000)(-0.3038) = 1.9411.
\end{align}
Jadi keluaran jaringan adalah $\Delta u = +1{,}94\%$. Sebagai pembanding, aksi yang
direkam dari kebijakan acuan pada kondisi tersebut adalah $+2{,}0\%$, sehingga selisih
prediksi hanya $0{,}06\%$ --- menunjukkan jaringan berhasil meniru kebijakan kendali.

\subsection{Akumulasi dan Penerapan}
\label{sub:akumulasi}
Keluaran jaringan berupa \emph{perubahan} \emph{duty cycle}. Nilai \emph{duty cycle}
aktual diperoleh dengan mengakumulasi dan membatasi pada rentang operasi
$[70\%,95\%]$:
\begin{equation}
u(k+1) = \mathrm{clip}\big(u(k) + \Delta u(k),\ 70,\ 95\big).
\label{eq:akumulasi}
\end{equation}
Untuk contoh di atas: $u(k+1)=\mathrm{clip}(82{,}00+1{,}94,\,70,\,95)=83{,}94\%$.
Sifat akumulatif ini memberi karakter integral pada pengendali.

Secara ringkas, seluruh model pengendali dapat ditulis sebagai satu fungsi:
\begin{equation}
\Delta u = \tanh\!\big(\mathbf{x'}\,W_1 + \mathbf{b}_1\big)\,W_2 + b_2,
\qquad \mathbf{x'} = (\mathbf{x}-\mu)/\sigma .
\label{eq:model-ringkas}
\end{equation}

\subsection{Pelatihan: Pembentukan Matriks Bobot}
\label{sub:pelatihan}
Nilai matriks pada Pers.~\eqref{eq:W1}--\eqref{eq:W2} \emph{tidak} ditetapkan manual,
melainkan diperoleh melalui pelatihan. Diberikan $N$ pasangan data
$(\mathbf{x}, \Delta u_{\text{target}})$, didefinisikan fungsi rugi (\emph{loss})
\emph{Mean Squared Error}:
\begin{equation}
L = \frac{1}{N}\sum_{i=1}^{N}\big(\Delta u^{(i)}_{\text{pred}} - \Delta u^{(i)}_{\text{target}}\big)^{2}.
\label{eq:loss}
\end{equation}
Gradien $\partial L/\partial W$ dihitung dengan algoritma \emph{backpropagation}
(aturan rantai dari keluaran ke masukan), lalu bobot diperbarui dengan
\emph{gradient descent}:
\begin{equation}
W \leftarrow W - \eta\,\frac{\partial L}{\partial W},
\label{eq:gd}
\end{equation}
dengan $\eta$ laju pembelajaran. Proses ini diulang hingga $L$ konvergen, sehingga
diperoleh nilai bobot akhir.

\subsection{Inferensi (\emph{Inference})}
\label{sub:inferensi}
\emph{Inferensi} adalah proses menjalankan jaringan \emph{maju saja}
(Pers.~\eqref{eq:scaling}, \eqref{eq:hidden}, \eqref{eq:output}, \eqref{eq:akumulasi})
menggunakan bobot yang telah dibekukan, tanpa menghitung gradien maupun memperbarui
bobot. Karena hanya melibatkan perkalian matriks, penjumlahan, dan fungsi tanh
(total 21 parameter), inferensi sangat ringan sehingga dapat dijalankan secara
\emph{real-time} pada mikrokontroler STM32 setiap selang waktu kendali. Dengan
demikian, \emph{pelatihan membentuk matriks bobot} (lewat \emph{backpropagation}),
sedangkan \emph{inferensi menggunakan matriks tersebut} (lewat propagasi maju).

 di Bab 2 (teori) / Bab 3 (perancangan).
% Catatan: desimal di dalam math memakai titik (.) agar spasi rapi; teks biasa
%          boleh memakai koma. Butuh paket amsmath (umumnya sudah dimuat).
% ============================================================================

\section{Dasar Jaringan Saraf Tiruan (Multilayer Perceptron)}
\label{sec:dasar-mlp}

Pengendali pada penelitian ini menggunakan jaringan saraf tiruan jenis
\emph{Multilayer Perceptron} (MLP), yaitu jaringan \emph{feedforward} di mana data
mengalir satu arah dari lapisan masukan, melalui lapisan tersembunyi, menuju lapisan
keluaran. MLP terdiri atas unit-unit hitung (neuron) yang saling terhubung melalui
\emph{bobot} (\emph{weight}) dan disertai \emph{bias}. Pada tiap neuron, masukan
berbobot dijumlahkan lalu dilewatkan ke sebuah \emph{fungsi aktivasi} non-linear.

Arsitektur yang dirancang adalah \textbf{3--4--1}, yaitu 3 neuron masukan, 4 neuron
pada satu lapisan tersembunyi, dan 1 neuron keluaran:
\begin{itemize}
  \item \textbf{Masukan} ($\mathbf{x}$): galat tekanan $e$, perubahan galat $\Delta e$,
        dan \emph{duty cycle} saat ini $u$.
  \item \textbf{Keluaran} ($y$): perubahan \emph{duty cycle} $\Delta u$ ($\Delta Duty$).
\end{itemize}

\subsection{Fungsi Aktivasi}
\label{sub:aktivasi}
Lapisan tersembunyi memakai fungsi aktivasi \emph{hyperbolic tangent} (tanh):
\begin{equation}
\tanh(z) = \frac{e^{z}-e^{-z}}{e^{z}+e^{-z}}, \qquad \tanh(z)\in(-1,1).
\label{eq:tanh}
\end{equation}
Fungsi non-linear ini penting: tanpa aktivasi non-linear, gabungan beberapa lapisan
linear akan setara dengan satu fungsi linear saja, sehingga jaringan tidak dapat
memodelkan hubungan melengkung. Lapisan keluaran memakai aktivasi linear (identitas)
agar nilai $\Delta u$ bebas bernilai positif, nol, maupun negatif.

\subsection{Propagasi Maju (\emph{Forward Propagation})}
\label{sub:forward}
Misalkan $\mathbf{x'}$ adalah vektor masukan (setelah dinormalisasi). Keluaran
lapisan tersembunyi dan lapisan keluaran dihitung sebagai:
\begin{align}
\mathbf{z} &= \mathbf{x'}\,W_1 + \mathbf{b}_1, &
\mathbf{h} &= \tanh(\mathbf{z}), \label{eq:hidden}\\
y &= \mathbf{h}\,W_2 + b_2, & \Delta u &= y. \label{eq:output}
\end{align}
dengan $W_1$ matriks bobot lapisan 1, $\mathbf{b}_1$ vektor bias lapisan 1, $W_2$
matriks bobot lapisan 2, dan $b_2$ bias keluaran. Persamaan~\eqref{eq:hidden} dan
\eqref{eq:output} adalah inti perhitungan MLP: perkalian matriks, penjumlahan bias,
dan fungsi aktivasi.

\subsection{Normalisasi Masukan}
\label{sub:scaling}
Karena ketiga masukan memiliki rentang sangat berbeda (misal $e\approx0{,}1$ tetapi
$u\approx80$), tiap masukan distandardisasi terlebih dahulu menggunakan rata-rata
$\mu_i$ dan simpangan baku $\sigma_i$ dari data latih:
\begin{equation}
x'_i = \frac{x_i - \mu_i}{\sigma_i}.
\label{eq:scaling}
\end{equation}

\section{Perancangan dan Model Matematis Pengendali MLP}
\label{sec:perancangan-mlp}

\subsection{Bobot Terlatih dan Dimensi Matriks}
\label{sub:bobot}
Setelah pelatihan, diperoleh nilai bobot dan bias yang tetap. Untuk arsitektur
3--4--1:
\begin{equation}
W_1 =
\begin{bmatrix}
-1.1576 & 1.7171 & -3.6581 & -16.9100\\
0.0077 & 0.1740 & -3.6078 & 7.4493\\
0.9206 & 0.7704 & -11.9742 & 10.4903
\end{bmatrix},\quad
\mathbf{b}_1 =
\begin{bmatrix}
0.5008\\ -0.3952\\ -22.8359\\ -7.2012
\end{bmatrix}^{\!\top}
\label{eq:W1}
\end{equation}
\begin{equation}
W_2 =
\begin{bmatrix}
-1.3646\\ 1.7042\\ 0.9956\\ -0.3038
\end{bmatrix},\qquad
b_2 = 1.2567 .
\label{eq:W2}
\end{equation}
Dimensi matriks ditentukan oleh aturan perkalian matriks. Karena 3 masukan menuju 4
neuron, $W_1$ berukuran $3\times4$ (tiap baris satu fitur, tiap kolom satu neuron).
Karena 4 neuron menuju 1 keluaran, $W_2$ berukuran $4\times1$. Dengan demikian:
\begin{equation}
\underbrace{(1\times3)}_{\mathbf{x'}}\cdot\underbrace{(3\times4)}_{W_1}=(1\times4),
\qquad
\underbrace{(1\times4)}_{\mathbf{h}}\cdot\underbrace{(4\times1)}_{W_2}=(1\times1),
\end{equation}
yaitu dimensi dalam harus sama ($3=3$ dan $4=4$). Jumlah parameter total adalah
$(3\times4+4)+(4\times1+1)=21$.

\subsection{Contoh Perhitungan Numerik (Data Nyata)}
\label{sub:contoh}
Diambil satu kondisi nyata: $P^{*}=0{,}300$~bar, $P=0{,}2408$~bar, sehingga
$e=0{,}0592$; $\Delta e=-0{,}0047$; dan $u=82{,}00\%$. Maka
$\mathbf{x}=[\,0.0592,\,-0.0047,\,82.00\,]$.

\textbf{Langkah 1 --- Normalisasi} (Pers.~\ref{eq:scaling}), dengan
$\mu=[-0.0159,\,0.0010,\,82.272]$ dan $\sigma=[0.1159,\,0.0616,\,7.1318]$:
\begin{equation}
\mathbf{x'} = [\,0.6475,\ -0.0931,\ -0.0382\,].
\end{equation}

\textbf{Langkah 2 --- Pra-aktivasi lapisan tersembunyi} (Pers.~\ref{eq:hidden}).
Sebagai contoh, neuron pertama (kolom ke-1 dari $W_1$):
\begin{equation}
z_1 = 0.5008 + (0.6475)(-1.1576) + (-0.0931)(0.0077) + (-0.0382)(0.9206) = -0.2845.
\end{equation}
Untuk keempat neuron diperoleh:
\begin{equation}
\mathbf{z} = [\,-0.2845,\ 0.6710,\ -24.4117,\ -19.2442\,].
\end{equation}

\textbf{Langkah 3 --- Aktivasi tanh} (Pers.~\ref{eq:tanh}):
\begin{equation}
\mathbf{h} = \tanh(\mathbf{z}) = [\,-0.2771,\ 0.5857,\ -1.0000,\ -1.0000\,].
\end{equation}
(Dua neuron terakhir jenuh pada $-1$ karena nilai $z$ sangat negatif.)

\textbf{Langkah 4 --- Lapisan keluaran} (Pers.~\ref{eq:output}):
\begin{align}
y &= 1.2567 + (-0.2771)(-1.3646) + (0.5857)(1.7042) \notag\\
  &\quad + (-1.0000)(0.9956) + (-1.0000)(-0.3038) = 1.9411.
\end{align}
Jadi keluaran jaringan adalah $\Delta u = +1{,}94\%$. Sebagai pembanding, aksi yang
direkam dari kebijakan acuan pada kondisi tersebut adalah $+2{,}0\%$, sehingga selisih
prediksi hanya $0{,}06\%$ --- menunjukkan jaringan berhasil meniru kebijakan kendali.

\subsection{Akumulasi dan Penerapan}
\label{sub:akumulasi}
Keluaran jaringan berupa \emph{perubahan} \emph{duty cycle}. Nilai \emph{duty cycle}
aktual diperoleh dengan mengakumulasi dan membatasi pada rentang operasi
$[70\%,95\%]$:
\begin{equation}
u(k+1) = \mathrm{clip}\big(u(k) + \Delta u(k),\ 70,\ 95\big).
\label{eq:akumulasi}
\end{equation}
Untuk contoh di atas: $u(k+1)=\mathrm{clip}(82{,}00+1{,}94,\,70,\,95)=83{,}94\%$.
Sifat akumulatif ini memberi karakter integral pada pengendali.

Secara ringkas, seluruh model pengendali dapat ditulis sebagai satu fungsi:
\begin{equation}
\Delta u = \tanh\!\big(\mathbf{x'}\,W_1 + \mathbf{b}_1\big)\,W_2 + b_2,
\qquad \mathbf{x'} = (\mathbf{x}-\mu)/\sigma .
\label{eq:model-ringkas}
\end{equation}

\subsection{Pelatihan: Pembentukan Matriks Bobot}
\label{sub:pelatihan}
Nilai matriks pada Pers.~\eqref{eq:W1}--\eqref{eq:W2} \emph{tidak} ditetapkan manual,
melainkan diperoleh melalui pelatihan. Diberikan $N$ pasangan data
$(\mathbf{x}, \Delta u_{\text{target}})$, didefinisikan fungsi rugi (\emph{loss})
\emph{Mean Squared Error}:
\begin{equation}
L = \frac{1}{N}\sum_{i=1}^{N}\big(\Delta u^{(i)}_{\text{pred}} - \Delta u^{(i)}_{\text{target}}\big)^{2}.
\label{eq:loss}
\end{equation}
Gradien $\partial L/\partial W$ dihitung dengan algoritma \emph{backpropagation}
(aturan rantai dari keluaran ke masukan), lalu bobot diperbarui dengan
\emph{gradient descent}:
\begin{equation}
W \leftarrow W - \eta\,\frac{\partial L}{\partial W},
\label{eq:gd}
\end{equation}
dengan $\eta$ laju pembelajaran. Proses ini diulang hingga $L$ konvergen, sehingga
diperoleh nilai bobot akhir.

\subsection{Inferensi (\emph{Inference})}
\label{sub:inferensi}
\emph{Inferensi} adalah proses menjalankan jaringan \emph{maju saja}
(Pers.~\eqref{eq:scaling}, \eqref{eq:hidden}, \eqref{eq:output}, \eqref{eq:akumulasi})
menggunakan bobot yang telah dibekukan, tanpa menghitung gradien maupun memperbarui
bobot. Karena hanya melibatkan perkalian matriks, penjumlahan, dan fungsi tanh
(total 21 parameter), inferensi sangat ringan sehingga dapat dijalankan secara
\emph{real-time} pada mikrokontroler STM32 setiap selang waktu kendali. Dengan
demikian, \emph{pelatihan membentuk matriks bobot} (lewat \emph{backpropagation}),
sedangkan \emph{inferensi menggunakan matriks tersebut} (lewat propagasi maju).

 di Bab 2 (teori) / Bab 3 (perancangan).
% Catatan: desimal di dalam math memakai titik (.) agar spasi rapi; teks biasa
%          boleh memakai koma. Butuh paket amsmath (umumnya sudah dimuat).
% ============================================================================

\section{Dasar Jaringan Saraf Tiruan (Multilayer Perceptron)}
\label{sec:dasar-mlp}

Pengendali pada penelitian ini menggunakan jaringan saraf tiruan jenis
\emph{Multilayer Perceptron} (MLP), yaitu jaringan \emph{feedforward} di mana data
mengalir satu arah dari lapisan masukan, melalui lapisan tersembunyi, menuju lapisan
keluaran. MLP terdiri atas unit-unit hitung (neuron) yang saling terhubung melalui
\emph{bobot} (\emph{weight}) dan disertai \emph{bias}. Pada tiap neuron, masukan
berbobot dijumlahkan lalu dilewatkan ke sebuah \emph{fungsi aktivasi} non-linear.

Arsitektur yang dirancang adalah \textbf{3--4--1}, yaitu 3 neuron masukan, 4 neuron
pada satu lapisan tersembunyi, dan 1 neuron keluaran:
\begin{itemize}
  \item \textbf{Masukan} ($\mathbf{x}$): galat tekanan $e$, perubahan galat $\Delta e$,
        dan \emph{duty cycle} saat ini $u$.
  \item \textbf{Keluaran} ($y$): perubahan \emph{duty cycle} $\Delta u$ ($\Delta Duty$).
\end{itemize}

\subsection{Fungsi Aktivasi}
\label{sub:aktivasi}
Lapisan tersembunyi memakai fungsi aktivasi \emph{hyperbolic tangent} (tanh):
\begin{equation}
\tanh(z) = \frac{e^{z}-e^{-z}}{e^{z}+e^{-z}}, \qquad \tanh(z)\in(-1,1).
\label{eq:tanh}
\end{equation}
Fungsi non-linear ini penting: tanpa aktivasi non-linear, gabungan beberapa lapisan
linear akan setara dengan satu fungsi linear saja, sehingga jaringan tidak dapat
memodelkan hubungan melengkung. Lapisan keluaran memakai aktivasi linear (identitas)
agar nilai $\Delta u$ bebas bernilai positif, nol, maupun negatif.

\subsection{Propagasi Maju (\emph{Forward Propagation})}
\label{sub:forward}
Misalkan $\mathbf{x'}$ adalah vektor masukan (setelah dinormalisasi). Keluaran
lapisan tersembunyi dan lapisan keluaran dihitung sebagai:
\begin{align}
\mathbf{z} &= \mathbf{x'}\,W_1 + \mathbf{b}_1, &
\mathbf{h} &= \tanh(\mathbf{z}), \label{eq:hidden}\\
y &= \mathbf{h}\,W_2 + b_2, & \Delta u &= y. \label{eq:output}
\end{align}
dengan $W_1$ matriks bobot lapisan 1, $\mathbf{b}_1$ vektor bias lapisan 1, $W_2$
matriks bobot lapisan 2, dan $b_2$ bias keluaran. Persamaan~\eqref{eq:hidden} dan
\eqref{eq:output} adalah inti perhitungan MLP: perkalian matriks, penjumlahan bias,
dan fungsi aktivasi.

\subsection{Normalisasi Masukan}
\label{sub:scaling}
Karena ketiga masukan memiliki rentang sangat berbeda (misal $e\approx0{,}1$ tetapi
$u\approx80$), tiap masukan distandardisasi terlebih dahulu menggunakan rata-rata
$\mu_i$ dan simpangan baku $\sigma_i$ dari data latih:
\begin{equation}
x'_i = \frac{x_i - \mu_i}{\sigma_i}.
\label{eq:scaling}
\end{equation}

\section{Perancangan dan Model Matematis Pengendali MLP}
\label{sec:perancangan-mlp}

\subsection{Bobot Terlatih dan Dimensi Matriks}
\label{sub:bobot}
Setelah pelatihan, diperoleh nilai bobot dan bias yang tetap. Untuk arsitektur
3--4--1:
\begin{equation}
W_1 =
\begin{bmatrix}
-1.1576 & 1.7171 & -3.6581 & -16.9100\\
0.0077 & 0.1740 & -3.6078 & 7.4493\\
0.9206 & 0.7704 & -11.9742 & 10.4903
\end{bmatrix},\quad
\mathbf{b}_1 =
\begin{bmatrix}
0.5008\\ -0.3952\\ -22.8359\\ -7.2012
\end{bmatrix}^{\!\top}
\label{eq:W1}
\end{equation}
\begin{equation}
W_2 =
\begin{bmatrix}
-1.3646\\ 1.7042\\ 0.9956\\ -0.3038
\end{bmatrix},\qquad
b_2 = 1.2567 .
\label{eq:W2}
\end{equation}
Dimensi matriks ditentukan oleh aturan perkalian matriks. Karena 3 masukan menuju 4
neuron, $W_1$ berukuran $3\times4$ (tiap baris satu fitur, tiap kolom satu neuron).
Karena 4 neuron menuju 1 keluaran, $W_2$ berukuran $4\times1$. Dengan demikian:
\begin{equation}
\underbrace{(1\times3)}_{\mathbf{x'}}\cdot\underbrace{(3\times4)}_{W_1}=(1\times4),
\qquad
\underbrace{(1\times4)}_{\mathbf{h}}\cdot\underbrace{(4\times1)}_{W_2}=(1\times1),
\end{equation}
yaitu dimensi dalam harus sama ($3=3$ dan $4=4$). Jumlah parameter total adalah
$(3\times4+4)+(4\times1+1)=21$.

\subsection{Contoh Perhitungan Numerik (Data Nyata)}
\label{sub:contoh}
Diambil satu kondisi nyata: $P^{*}=0{,}300$~bar, $P=0{,}2408$~bar, sehingga
$e=0{,}0592$; $\Delta e=-0{,}0047$; dan $u=82{,}00\%$. Maka
$\mathbf{x}=[\,0.0592,\,-0.0047,\,82.00\,]$.

\textbf{Langkah 1 --- Normalisasi} (Pers.~\ref{eq:scaling}), dengan
$\mu=[-0.0159,\,0.0010,\,82.272]$ dan $\sigma=[0.1159,\,0.0616,\,7.1318]$:
\begin{equation}
\mathbf{x'} = [\,0.6475,\ -0.0931,\ -0.0382\,].
\end{equation}

\textbf{Langkah 2 --- Pra-aktivasi lapisan tersembunyi} (Pers.~\ref{eq:hidden}).
Sebagai contoh, neuron pertama (kolom ke-1 dari $W_1$):
\begin{equation}
z_1 = 0.5008 + (0.6475)(-1.1576) + (-0.0931)(0.0077) + (-0.0382)(0.9206) = -0.2845.
\end{equation}
Untuk keempat neuron diperoleh:
\begin{equation}
\mathbf{z} = [\,-0.2845,\ 0.6710,\ -24.4117,\ -19.2442\,].
\end{equation}

\textbf{Langkah 3 --- Aktivasi tanh} (Pers.~\ref{eq:tanh}):
\begin{equation}
\mathbf{h} = \tanh(\mathbf{z}) = [\,-0.2771,\ 0.5857,\ -1.0000,\ -1.0000\,].
\end{equation}
(Dua neuron terakhir jenuh pada $-1$ karena nilai $z$ sangat negatif.)

\textbf{Langkah 4 --- Lapisan keluaran} (Pers.~\ref{eq:output}):
\begin{align}
y &= 1.2567 + (-0.2771)(-1.3646) + (0.5857)(1.7042) \notag\\
  &\quad + (-1.0000)(0.9956) + (-1.0000)(-0.3038) = 1.9411.
\end{align}
Jadi keluaran jaringan adalah $\Delta u = +1{,}94\%$. Sebagai pembanding, aksi yang
direkam dari kebijakan acuan pada kondisi tersebut adalah $+2{,}0\%$, sehingga selisih
prediksi hanya $0{,}06\%$ --- menunjukkan jaringan berhasil meniru kebijakan kendali.

\subsection{Akumulasi dan Penerapan}
\label{sub:akumulasi}
Keluaran jaringan berupa \emph{perubahan} \emph{duty cycle}. Nilai \emph{duty cycle}
aktual diperoleh dengan mengakumulasi dan membatasi pada rentang operasi
$[70\%,95\%]$:
\begin{equation}
u(k+1) = \mathrm{clip}\big(u(k) + \Delta u(k),\ 70,\ 95\big).
\label{eq:akumulasi}
\end{equation}
Untuk contoh di atas: $u(k+1)=\mathrm{clip}(82{,}00+1{,}94,\,70,\,95)=83{,}94\%$.
Sifat akumulatif ini memberi karakter integral pada pengendali.

Secara ringkas, seluruh model pengendali dapat ditulis sebagai satu fungsi:
\begin{equation}
\Delta u = \tanh\!\big(\mathbf{x'}\,W_1 + \mathbf{b}_1\big)\,W_2 + b_2,
\qquad \mathbf{x'} = (\mathbf{x}-\mu)/\sigma .
\label{eq:model-ringkas}
\end{equation}

\subsection{Pelatihan: Pembentukan Matriks Bobot}
\label{sub:pelatihan}
Nilai matriks pada Pers.~\eqref{eq:W1}--\eqref{eq:W2} \emph{tidak} ditetapkan manual,
melainkan diperoleh melalui pelatihan. Diberikan $N$ pasangan data
$(\mathbf{x}, \Delta u_{\text{target}})$, didefinisikan fungsi rugi (\emph{loss})
\emph{Mean Squared Error}:
\begin{equation}
L = \frac{1}{N}\sum_{i=1}^{N}\big(\Delta u^{(i)}_{\text{pred}} - \Delta u^{(i)}_{\text{target}}\big)^{2}.
\label{eq:loss}
\end{equation}
Gradien $\partial L/\partial W$ dihitung dengan algoritma \emph{backpropagation}
(aturan rantai dari keluaran ke masukan), lalu bobot diperbarui dengan
\emph{gradient descent}:
\begin{equation}
W \leftarrow W - \eta\,\frac{\partial L}{\partial W},
\label{eq:gd}
\end{equation}
dengan $\eta$ laju pembelajaran. Proses ini diulang hingga $L$ konvergen, sehingga
diperoleh nilai bobot akhir.

\subsection{Inferensi (\emph{Inference})}
\label{sub:inferensi}
\emph{Inferensi} adalah proses menjalankan jaringan \emph{maju saja}
(Pers.~\eqref{eq:scaling}, \eqref{eq:hidden}, \eqref{eq:output}, \eqref{eq:akumulasi})
menggunakan bobot yang telah dibekukan, tanpa menghitung gradien maupun memperbarui
bobot. Karena hanya melibatkan perkalian matriks, penjumlahan, dan fungsi tanh
(total 21 parameter), inferensi sangat ringan sehingga dapat dijalankan secara
\emph{real-time} pada mikrokontroler STM32 setiap selang waktu kendali. Dengan
demikian, \emph{pelatihan membentuk matriks bobot} (lewat \emph{backpropagation}),
sedangkan \emph{inferensi menggunakan matriks tersebut} (lewat propagasi maju).

 di Bab 2 (teori) / Bab 3 (perancangan).
% Catatan: desimal di dalam math memakai titik (.) agar spasi rapi; teks biasa
%          boleh memakai koma. Butuh paket amsmath (umumnya sudah dimuat).
% ============================================================================

\section{Dasar Jaringan Saraf Tiruan (Multilayer Perceptron)}
\label{sec:dasar-mlp}

Pengendali pada penelitian ini menggunakan jaringan saraf tiruan jenis
\emph{Multilayer Perceptron} (MLP), yaitu jaringan \emph{feedforward} di mana data
mengalir satu arah dari lapisan masukan, melalui lapisan tersembunyi, menuju lapisan
keluaran. MLP terdiri atas unit-unit hitung (neuron) yang saling terhubung melalui
\emph{bobot} (\emph{weight}) dan disertai \emph{bias}. Pada tiap neuron, masukan
berbobot dijumlahkan lalu dilewatkan ke sebuah \emph{fungsi aktivasi} non-linear.

Arsitektur yang dirancang adalah \textbf{3--4--1}, yaitu 3 neuron masukan, 4 neuron
pada satu lapisan tersembunyi, dan 1 neuron keluaran:
\begin{itemize}
  \item \textbf{Masukan} ($\mathbf{x}$): galat tekanan $e$, perubahan galat $\Delta e$,
        dan \emph{duty cycle} saat ini $u$.
  \item \textbf{Keluaran} ($y$): perubahan \emph{duty cycle} $\Delta u$ ($\Delta Duty$).
\end{itemize}

\subsection{Fungsi Aktivasi}
\label{sub:aktivasi}
Lapisan tersembunyi memakai fungsi aktivasi \emph{hyperbolic tangent} (tanh):
\begin{equation}
\tanh(z) = \frac{e^{z}-e^{-z}}{e^{z}+e^{-z}}, \qquad \tanh(z)\in(-1,1).
\label{eq:tanh}
\end{equation}
Fungsi non-linear ini penting: tanpa aktivasi non-linear, gabungan beberapa lapisan
linear akan setara dengan satu fungsi linear saja, sehingga jaringan tidak dapat
memodelkan hubungan melengkung. Lapisan keluaran memakai aktivasi linear (identitas)
agar nilai $\Delta u$ bebas bernilai positif, nol, maupun negatif.

\subsection{Propagasi Maju (\emph{Forward Propagation})}
\label{sub:forward}
Misalkan $\mathbf{x'}$ adalah vektor masukan (setelah dinormalisasi). Keluaran
lapisan tersembunyi dan lapisan keluaran dihitung sebagai:
\begin{align}
\mathbf{z} &= \mathbf{x'}\,W_1 + \mathbf{b}_1, &
\mathbf{h} &= \tanh(\mathbf{z}), \label{eq:hidden}\\
y &= \mathbf{h}\,W_2 + b_2, & \Delta u &= y. \label{eq:output}
\end{align}
dengan $W_1$ matriks bobot lapisan 1, $\mathbf{b}_1$ vektor bias lapisan 1, $W_2$
matriks bobot lapisan 2, dan $b_2$ bias keluaran. Persamaan~\eqref{eq:hidden} dan
\eqref{eq:output} adalah inti perhitungan MLP: perkalian matriks, penjumlahan bias,
dan fungsi aktivasi.

\subsection{Normalisasi Masukan}
\label{sub:scaling}
Karena ketiga masukan memiliki rentang sangat berbeda (misal $e\approx0{,}1$ tetapi
$u\approx80$), tiap masukan distandardisasi terlebih dahulu menggunakan rata-rata
$\mu_i$ dan simpangan baku $\sigma_i$ dari data latih:
\begin{equation}
x'_i = \frac{x_i - \mu_i}{\sigma_i}.
\label{eq:scaling}
\end{equation}

\section{Perancangan dan Model Matematis Pengendali MLP}
\label{sec:perancangan-mlp}

\subsection{Bobot Terlatih dan Dimensi Matriks}
\label{sub:bobot}
Setelah pelatihan, diperoleh nilai bobot dan bias yang tetap. Untuk arsitektur
3--4--1:
\begin{equation}
W_1 =
\begin{bmatrix}
-1.1576 & 1.7171 & -3.6581 & -16.9100\\
0.0077 & 0.1740 & -3.6078 & 7.4493\\
0.9206 & 0.7704 & -11.9742 & 10.4903
\end{bmatrix},\quad
\mathbf{b}_1 =
\begin{bmatrix}
0.5008\\ -0.3952\\ -22.8359\\ -7.2012
\end{bmatrix}^{\!\top}
\label{eq:W1}
\end{equation}
\begin{equation}
W_2 =
\begin{bmatrix}
-1.3646\\ 1.7042\\ 0.9956\\ -0.3038
\end{bmatrix},\qquad
b_2 = 1.2567 .
\label{eq:W2}
\end{equation}
Dimensi matriks ditentukan oleh aturan perkalian matriks. Karena 3 masukan menuju 4
neuron, $W_1$ berukuran $3\times4$ (tiap baris satu fitur, tiap kolom satu neuron).
Karena 4 neuron menuju 1 keluaran, $W_2$ berukuran $4\times1$. Dengan demikian:
\begin{equation}
\underbrace{(1\times3)}_{\mathbf{x'}}\cdot\underbrace{(3\times4)}_{W_1}=(1\times4),
\qquad
\underbrace{(1\times4)}_{\mathbf{h}}\cdot\underbrace{(4\times1)}_{W_2}=(1\times1),
\end{equation}
yaitu dimensi dalam harus sama ($3=3$ dan $4=4$). Jumlah parameter total adalah
$(3\times4+4)+(4\times1+1)=21$.

\subsection{Contoh Perhitungan Numerik (Data Nyata)}
\label{sub:contoh}
Diambil satu kondisi nyata: $P^{*}=0{,}300$~bar, $P=0{,}2408$~bar, sehingga
$e=0{,}0592$; $\Delta e=-0{,}0047$; dan $u=82{,}00\%$. Maka
$\mathbf{x}=[\,0.0592,\,-0.0047,\,82.00\,]$.

\textbf{Langkah 1 --- Normalisasi} (Pers.~\ref{eq:scaling}), dengan
$\mu=[-0.0159,\,0.0010,\,82.272]$ dan $\sigma=[0.1159,\,0.0616,\,7.1318]$:
\begin{equation}
\mathbf{x'} = [\,0.6475,\ -0.0931,\ -0.0382\,].
\end{equation}

\textbf{Langkah 2 --- Pra-aktivasi lapisan tersembunyi} (Pers.~\ref{eq:hidden}).
Sebagai contoh, neuron pertama (kolom ke-1 dari $W_1$):
\begin{equation}
z_1 = 0.5008 + (0.6475)(-1.1576) + (-0.0931)(0.0077) + (-0.0382)(0.9206) = -0.2845.
\end{equation}
Untuk keempat neuron diperoleh:
\begin{equation}
\mathbf{z} = [\,-0.2845,\ 0.6710,\ -24.4117,\ -19.2442\,].
\end{equation}

\textbf{Langkah 3 --- Aktivasi tanh} (Pers.~\ref{eq:tanh}):
\begin{equation}
\mathbf{h} = \tanh(\mathbf{z}) = [\,-0.2771,\ 0.5857,\ -1.0000,\ -1.0000\,].
\end{equation}
(Dua neuron terakhir jenuh pada $-1$ karena nilai $z$ sangat negatif.)

\textbf{Langkah 4 --- Lapisan keluaran} (Pers.~\ref{eq:output}):
\begin{align}
y &= 1.2567 + (-0.2771)(-1.3646) + (0.5857)(1.7042) \notag\\
  &\quad + (-1.0000)(0.9956) + (-1.0000)(-0.3038) = 1.9411.
\end{align}
Jadi keluaran jaringan adalah $\Delta u = +1{,}94\%$. Sebagai pembanding, aksi yang
direkam dari kebijakan acuan pada kondisi tersebut adalah $+2{,}0\%$, sehingga selisih
prediksi hanya $0{,}06\%$ --- menunjukkan jaringan berhasil meniru kebijakan kendali.

\subsection{Akumulasi dan Penerapan}
\label{sub:akumulasi}
Keluaran jaringan berupa \emph{perubahan} \emph{duty cycle}. Nilai \emph{duty cycle}
aktual diperoleh dengan mengakumulasi dan membatasi pada rentang operasi
$[70\%,95\%]$:
\begin{equation}
u(k+1) = \mathrm{clip}\big(u(k) + \Delta u(k),\ 70,\ 95\big).
\label{eq:akumulasi}
\end{equation}
Untuk contoh di atas: $u(k+1)=\mathrm{clip}(82{,}00+1{,}94,\,70,\,95)=83{,}94\%$.
Sifat akumulatif ini memberi karakter integral pada pengendali.

Secara ringkas, seluruh model pengendali dapat ditulis sebagai satu fungsi:
\begin{equation}
\Delta u = \tanh\!\big(\mathbf{x'}\,W_1 + \mathbf{b}_1\big)\,W_2 + b_2,
\qquad \mathbf{x'} = (\mathbf{x}-\mu)/\sigma .
\label{eq:model-ringkas}
\end{equation}

\subsection{Pelatihan: Pembentukan Matriks Bobot}
\label{sub:pelatihan}
Nilai matriks pada Pers.~\eqref{eq:W1}--\eqref{eq:W2} \emph{tidak} ditetapkan manual,
melainkan diperoleh melalui pelatihan. Diberikan $N$ pasangan data
$(\mathbf{x}, \Delta u_{\text{target}})$, didefinisikan fungsi rugi (\emph{loss})
\emph{Mean Squared Error}:
\begin{equation}
L = \frac{1}{N}\sum_{i=1}^{N}\big(\Delta u^{(i)}_{\text{pred}} - \Delta u^{(i)}_{\text{target}}\big)^{2}.
\label{eq:loss}
\end{equation}
Gradien $\partial L/\partial W$ dihitung dengan algoritma \emph{backpropagation}
(aturan rantai dari keluaran ke masukan), lalu bobot diperbarui dengan
\emph{gradient descent}:
\begin{equation}
W \leftarrow W - \eta\,\frac{\partial L}{\partial W},
\label{eq:gd}
\end{equation}
dengan $\eta$ laju pembelajaran. Proses ini diulang hingga $L$ konvergen, sehingga
diperoleh nilai bobot akhir.

\subsection{Inferensi (\emph{Inference})}
\label{sub:inferensi}
\emph{Inferensi} adalah proses menjalankan jaringan \emph{maju saja}
(Pers.~\eqref{eq:scaling}, \eqref{eq:hidden}, \eqref{eq:output}, \eqref{eq:akumulasi})
menggunakan bobot yang telah dibekukan, tanpa menghitung gradien maupun memperbarui
bobot. Karena hanya melibatkan perkalian matriks, penjumlahan, dan fungsi tanh
(total 21 parameter), inferensi sangat ringan sehingga dapat dijalankan secara
\emph{real-time} pada mikrokontroler STM32 setiap selang waktu kendali. Dengan
demikian, \emph{pelatihan membentuk matriks bobot} (lewat \emph{backpropagation}),
sedangkan \emph{inferensi menggunakan matriks tersebut} (lewat propagasi maju).

